% !TeX root = ../main.tex
\chapter*{Introducción}
\phantomsection
\addcontentsline{toc}{chapter}{Introducción}

\section*{Fundamentación}

La física acústica aporta modelos para describir cómo se generan, se propagan
y se registran las perturbaciones sonoras. En Fonoaudiología, esos modelos
permiten interpretar con precisión magnitudes, señales e instrumentos que
aparecen en el estudio de la audición, la voz, los ambientes y los
dispositivos de rehabilitación. Esta base no reemplaza el razonamiento
perceptual ni clínico: permite distinguir qué se mide físicamente, qué
experimenta una persona y qué conclusión profesional puede sostenerse a
partir de cada dato.

El libro sigue las diez unidades del programa de la materia. Las
Unidades~\ref{chap:unidad-1} a~\ref{chap:unidad-5} construyen la base
matemática, mecánica, ondulatoria y acústica, incluida la descripción
frecuencial de señales y sistemas. Las Unidades~\ref{chap:unidad-6} a
\ref{chap:unidad-8} estudian el mecanismo auditivo periférico, los fenómenos
psicoacústicos y la relación entre alteraciones, estudios y tecnologías de
rehabilitación. Las Unidades~\ref{chap:unidad-9} y~\ref{chap:unidad-10}
integran esos contenidos en el análisis de la propagación, los recintos, el
ruido y el enmascaramiento.

\section*{Objetivos generales y específicos}

\subsection*{Objetivos generales}

\begin{itemize}
    \setlength\itemsep{0.3em}
    \item Construir una perspectiva integrada de la acústica física, la
    percepción auditiva y sus aplicaciones fonoaudiológicas.
    \item Interpretar modelos, mediciones y representaciones sin confundir
    magnitudes físicas, atributos perceptuales y conclusiones clínicas.
\end{itemize}

\subsection*{Objetivos específicos}

\begin{itemize}
    \item Delimitar el campo de la acústica y reconocer su relación con la
    formación y la práctica fonoaudiológicas.
    \item Operar con magnitudes, unidades, gráficos y relaciones matemáticas
    básicas empleadas para describir fenómenos sonoros.
    \item Explicar la generación y la propagación del sonido mediante modelos
    mecánicos y ondulatorios, indicando sus hipótesis y límites.
    \item Diferenciar presión acústica, intensidad, potencia y sus niveles, así
    como frecuencia, altura tonal, nivel de presión sonora y sonoridad.
    \item Analizar señales y sistemas en los dominios temporal y frecuencial e
    interpretar espectros, filtros y respuestas en frecuencia.
    \item Describir las transformaciones físicas y fisiológicas del mecanismo
    auditivo periférico y relacionarlas con fenómenos psicoacústicos.
    \item Distinguir exposición, alteración, síntoma y resultado de un estudio,
    sin atribuir a una medición aislada un valor diagnóstico que no posee.
    \item Reconocer los principios físicos de estudios auditivos y
    tecnologías de rehabilitación, junto con sus alcances y limitaciones.
    \item Caracterizar la propagación en ambientes, el ruido y el
    enmascaramiento mediante descriptores físicos y criterios de medición
    explícitos.
\end{itemize}

\section*{Cómo recorrer el libro}

Cada unidad conserva un problema central propio y explicita los conocimientos
que recibe de las anteriores. El desarrollo parte del fenómeno y su
interpretación antes de introducir ecuaciones; los ejemplos resueltos muestran
el procedimiento y los ejercicios avanzan desde el reconocimiento conceptual
hacia la integración. Cuando un tema ya fue definido, las unidades posteriores
lo recuperan mediante una referencia y concentran el desarrollo en la nueva
aplicación.

%%%%%%%%%%%%%%%%%%%%%%%%%%%%%%%%%%%%%%%%%%%%%%%%%%%%%%%%%%%%%%%%%%%%%%%%%%%%%%%%%%%%%%%%%%%%%%%%%%
